%======================================================================
% TÓM TẮT ĐỒ ÁN
%======================================================================
\chapter*{Tóm tắt đồ án}
\addcontentsline{toc}{chapter}{Tóm tắt đồ án}

Đồ án nghiên cứu và triển khai mô hình bảo mật Zero Trust cho hệ thống tuyển dụng Job7189 xây dựng theo kiến trúc microservices trên Kubernetes. Hệ thống còn tồn tại mức tin cậy ngầm định trong mạng nội bộ, thông tin xác thực phụ thuộc vào cấu hình tĩnh và nguy cơ truy cập ngang khi một container bị xâm phạm.

Để giải quyết vấn đề này, đồ án khảo sát hiện trạng, xây dựng bản đồ luồng mạng, sau đó thiết kế và triển khai các nhóm chức năng: định danh, phân đoạn mạng, kiểm soát truy cập, quản lý bí mật động, giám sát thời gian chạy và quan sát nhật ký. Môi trường thử nghiệm sử dụng cụm Kubernetes nhiều nút trên các máy chủ ảo, kết nối qua Tailscale, cùng các công cụ như Cilium, Hubble, Kong Gateway, Keycloak, SPIRE, Vault, Tetragon, Trivy Operator, OPA Gatekeeper, Prometheus, Grafana và EFK.

Kết quả kiểm chứng cho thấy hệ thống chặn được nhiều hành vi như truy cập trái phép API, quét mạng nội bộ, kết nối cơ sở dữ liệu không hợp lệ, chạy công cụ bị cấm, dùng image chưa được phê duyệt và gửi dữ liệu ra ngoài. Đồ án chứng minh mô hình đề xuất có tính khả thi trong môi trường phòng thí nghiệm và có thể mở rộng theo hướng tự động hóa phản hồi, tích hợp thêm nguồn tình báo đe dọa và hoàn thiện cơ chế chấm điểm rủi ro. Qua quá trình thực hiện, sinh viên củng cố kiến thức về Kubernetes, bảo mật hệ thống phân tán, vận hành hạ tầng và kiểm chứng an ninh thực nghiệm.

\vspace{0.5cm}
\noindent\textbf{Từ khóa:} Zero Trust Architecture, NIST SP 800-207, Microservices Security, Kubernetes, Microsegmentation, Dynamic Secrets.
